\newpage
\section{Materiais e métodos}
\label{sc:materiais_e_metodos}

O desenvolvimento do Dança~AI no âmbito do TCC02 compreende quatro fases sequenciais,
alinhadas aos objetivos específicos 3 a 7 definidos na Seção~\ref{sec:objetivos}:
(i)~implementação do \textit{pipeline} de visão computacional e dos módulos de transferência
de peso e postura; (ii)~implementação dos módulos de ritmo e identificação de movimentos;
(iii)~integração dos quatro módulos em interface unificada e calibração; e
(iv)~avaliação do protótipo e redação da monografia. Esta seção descreve os materiais,
ferramentas e procedimentos metodológicos adotados.

\subsection{Hardware e Dispositivos-Alvo}
\label{sub:hardware}

Os testes do Dança~AI serão conduzidos em dispositivos Android de segmento intermediário
(\textit{mid-range}), com especificações mínimas de: processador octa-core a ao menos
1{,}8~GHz, GPU compatível com OpenGL ES~3.1 --- requisito do \textit{delegate} GPU do
TensorFlow Lite~\cite{lee2019ondevice} ---, 4~GB de RAM e câmera traseira com resolução
mínima de 1080p a 30~FPS. Esse perfil corresponde a dispositivos populares na faixa de preço
acessível ao público-alvo, garantindo que os resultados de desempenho obtidos sejam
representativos de condições reais de uso.

Os experimentos de latência serão realizados em ao menos dois dispositivos de especificações
distintas para verificar a robustez do orçamento de 100~ms estabelecido na
Seção~\ref{subsub:otimizacao} em hardware variado. O modelo BlazePose \textit{lite},
adotado no projeto, foi validado pelos seus autores a mais de 30~FPS no Pixel~2 --- dispositivo
lançado em 2017 com hardware significativamente inferior ao dos \textit{mid-range} atuais ---,
fornecendo margem de segurança confortável para execução no perfil-alvo~\cite{bazarevsky2020blazeposeondevicerealtimebody}.

\subsection{Stack Tecnológico}
\label{sub:stack}

O aplicativo é desenvolvido integralmente em Kotlin sobre a plataforma Android, com suporte
ao SDK mínimo 24 (Android 7.0 Nougat), representando mais de 94\% dos dispositivos ativos na
plataforma. A adoção do Kotlin é fundamentada nas garantias de segurança de tipos nulos em
tempo de compilação e no suporte nativo a corrotinas, que simplificam o gerenciamento do
\textit{pipeline} assíncrono de captura, inferência e análise de vídeo sem o \textit{overhead}
de \textit{callbacks} aninhados~\cite{jemerov2017kotlin}.

A captura de vídeo é gerenciada pela biblioteca CameraX do Jetpack Android, que fornece API
unificada e independente de fabricante para acesso à câmera. O caso de uso
\texttt{ImageAnalysis} entrega quadros no formato \texttt{ImageProxy} a um executor
configurável com controle de \textit{backpressure}, descartando quadros excedentes quando o
processador não acompanha a taxa de captura --- comportamento essencial para preservar a
responsividade da interface sob carga computacional elevada.

A estimativa de pose é realizada pelo MediaPipe Pose Landmarker~\cite{lugaresi2019mediapipe}
sobre o modelo BlazePose \textit{lite}~\cite{bazarevsky2020blazeposeondevicerealtimebody},
integrado ao Android via dependência Maven
\texttt{com.google.mediapipe:tasks-vision:0.10.14}. O modelo opera em modo
\texttt{LIVE\_STREAM}, processando cada quadro de forma assíncrona e retornando os
33~\textit{landmarks} tridimensionais com \textit{visibility scores} por \textit{callback}.
A execução é acelerada pelo \textit{delegate} GPU, com \textit{fallback} automático para CPU
em dispositivos sem suporte a OpenGL ES~3.1, eliminando qualquer dependência de hardware
especializado~\cite{lee2019ondevice}.

O módulo de comparação de movimentos implementa o algoritmo \textit{Dynamic Time Warping}
(DTW) sobre as séries temporais de ângulos articulares extraídas pelo
\texttt{AngleCalculator}, calculados \textit{frame} a \textit{frame} conforme a
Equação~\ref{eq:angulo}. A segmentação temporal --- isto é, a delimitação do início e do
fim de cada passo para montagem da sequência comparada --- é realizada pelo próprio módulo
de transferência de peso: um ciclo completo de transferência (ex.: esquerdo → direito →
esquerdo) sinaliza a conclusão de uma unidade de movimento e dispara a comparação DTW.
A implementação é realizada em Kotlin puro, sem dependências externas, com restrição de
janela de Sakoe-Chiba para limitar o custo computacional e tornar o alinhamento compatível
com execução em tempo real~\cite{dynamic_dance_warping}. Os \textit{templates} de
referência são gravados a partir de execuções de um instrutor qualificado e armazenados
localmente no dispositivo como séries de \texttt{BodyAngles} serializadas, eliminando
qualquer dependência de conectividade de rede em ambiente de prática.

Os módulos de postura e de identificação de movimentos utilizam modelos de classificação
no formato TensorFlow Lite (TFLite), executados via \texttt{org.tensorflow:tensorflow-lite},
com quantização em inteiros de 8~bits para reduzir o tamanho do modelo e o consumo de
memória em dispositivos \textit{mid-range}~\cite{bursa2023tflite}. Os modelos são treinados em Python com TensorFlow/Keras~\cite{abadi2016tensorflow}
em ambiente Google Colab, a partir de coordenadas $xyz$ dos 33~\textit{landmarks}
extraídas pelo MediaPipe, sem retenção dos vídeos brutos, preservando a privacidade
dos participantes e eliminando variações irrelevantes de iluminação, vestuário e fundo.
A conversão para o formato \texttt{.tflite} é realizada via
\texttt{tf.lite.TFLiteConverter}~\cite{lee2019ondevice}.
O classificador de postura opera sobre valores instantâneos de um único quadro;
o classificador de movimentos opera sobre séries temporais de quadros consecutivos. O dataset será coletado com múltiplos dançarinos para
ampliar a variabilidade inter-sujeito, com mínimo de 100 amostras por classe após
\textit{data augmentation} por espelhamento lateral e variação de velocidade; a anotação
das amostras será realizada pelos próprios autores, que possuem formação em dança de
salão. O classificador de postura é multi-label, avaliando cinco desvios de forma
independente --- ombro esquerdo caído, ombro direito caído, ombros encurvados,
inclinação do tronco para esquerda e inclinação do tronco para direita ---
podendo detectar múltiplos desvios simultaneamente; ausência de todos os
desvios corresponde à postura correta; o classificador de movimentos cobre
exclusivamente a identificação do passo em execução, com a avaliação de
qualidade delegada ao módulo DTW. Os modelos convertidos
para o formato \texttt{.tflite} são empacotados como \textit{assets} no APK.

O módulo de ritmo utiliza a API \texttt{AudioTrack} do Android para geração de sinais
sonoros com baixa latência, emitindo pulsações no BPM configurado pelo praticante. A
validação rítmica compara o intervalo inter-batida (IBI) do metrônomo com a cadência das
transferências de peso detectadas pelo respectivo módulo, gerando \textit{feedback} visual
sobre a conformidade rítmica em tempo real~\cite{rose2019metronome}.

A interface do usuário é construída com Jetpack Compose e \textit{Material~3}. A tela de
configuração de sessão inclui um guia de enquadramento que fornece feedback visual em tempo
real sobre a distância, ângulo e altura do dispositivo em relação ao praticante, orientando
o posicionamento antes do início da sessão; o uso de suporte físico é recomendado para
maior estabilidade, mas não é requisito do sistema. O \textit{overlay} de
\textit{feedback} visual é renderizado por uma \texttt{View} customizada
(\texttt{OverlayView}) sobreposta ao \textit{preview} da câmera, com indicadores visuais por módulo; no módulo de identificação de movimentos, a
codificação cromática por articulação é calculada a partir da distância DTW do
segmento corrente em relação ao \textit{template}. O controle de versão é realizado via Git com repositório remoto no GitHub,
adotando fluxo de ramificação por funcionalidade (\textit{feature branching}) que permite o
desenvolvimento paralelo dos módulos de pose, DTW e interface.

\subsection{Metodologia de Desenvolvimento}
\label{sub:metodologia}

O desenvolvimento segue abordagem iterativa orientada ao produto mínimo viável (MVP),
estruturada em quatro fases correspondentes aos objetivos específicos 3 a 7 do TCC02. Cada
fase produz um artefato verificável que serve de entrada para a fase seguinte, permitindo
detecção precoce de problemas de integração.

Na \textbf{Fase~1} são implementados o \textit{pipeline} de visão computacional e os
primeiros dois módulos de análise: captura via CameraX, inferência com MediaPipe Pose
Landmarker e renderização do \textit{overlay} de esqueleto; módulo de transferência de peso
com classificação esquerdo/direito/neutro e \textit{feedback} visual; e módulo de postura
com regras geométricas sobre os \textit{landmarks} da parte superior do corpo e modelo TFLite
embarcado treinado sobre coordenadas $xyz$ de dançarinos qualificados~\cite{roggio2024pose}. O critério de
conclusão é execução estável a ao menos 20~FPS com latência ponta a ponta inferior a 100~ms,
aferida via \texttt{SystemClock.uptimeMillis()}.

Na \textbf{Fase~2} são implementados o módulo de ritmo --- metrônomo configurável por BPM,
\textit{feedback} auditivo via \texttt{AudioTrack} e validação pela comparação entre o IBI
e a cadência de transferências de peso detectada em tempo real --- e o módulo de identificação
de movimentos, com treinamento e \textit{deploy} do classificador TFLite para os seis passos
do forró~\cite{skeleton_dtw_sport} e implementação dos modos Livre e Desafio. Os
\textit{templates} DTW de referência são gravados com instrutor qualificado de forró
universitário e armazenados localmente no dispositivo.

Na \textbf{Fase~3} os quatro módulos são integrados em interface unificada construída com
Jetpack Compose e \textit{Material~3}, e os limiares de aceitação dos classificadores e do
metrônomo são calibrados em sessões piloto com praticantes. O critério de conclusão é a
confirmação de que o \textit{feedback} é percebido como responsivo e coerente com o
movimento, dentro do limiar neurológico de 200~ms documentado na literatura de aprendizagem
motora~\cite{plos_feedback_delay}.

Na \textbf{Fase~4} os testes com usuários são conduzidos e a monografia é redigida,
conforme detalhado na próxima subseção.

\subsection{Protocolo de Avaliação}
\label{sub:avaliacao}

A avaliação do protótipo contempla quatro dimensões complementares, correspondentes ao
objetivo específico~7: desempenho computacional do \textit{pipeline}, acurácia da estimativa
de pose em condições reais, acurácia dos modelos de classificação embarcados e eficácia
pedagógica do sistema de \textit{feedback}.

\subsubsection{Desempenho Computacional}
A latência do \textit{pipeline} completo será medida quadro a quadro por estágio
(captura, inferência MediaPipe, inferência dos classificadores TFLite, módulo de ritmo e
renderização). As métricas
reportadas incluirão mediana e percentil~95 da latência de ponta a ponta, taxa de quadros por
segundo, percentual de quadros descartados por \textit{backpressure} e consumo de memória RAM
durante sessões de 5~minutos. Os testes serão conduzidos em dois dispositivos com
especificações contrastantes para verificar a robustez do orçamento de latência em hardware
variado.

\subsubsection{Acurácia da Estimativa de Pose}
A validade das medições angulares extraídas pelo MediaPipe Pose Landmarker é assumida
com base nos limites de concordância reportados por Latreche et
al.~\cite{LATRECHE2023112826}, que compararam os ângulos calculados a partir dos
\textit{landmarks} do MediaPipe com goniometria clínica e obtiveram valores dentro de
margens clinicamente aceitáveis para as articulações de joelhos, quadris e ombros ---
domínio com requisitos de precisão comparáveis aos do Dança~AI.

\subsubsection{Acurácia dos Modelos de Classificação}
A acurácia dos classificadores TFLite de postura e de identificação de movimentos será
avaliada em conjunto de teste com rótulos conhecidos: para o classificador de postura,
os participantes serão fotografados adotando deliberadamente cada uma das cinco categorias
de desvio e a postura correta, de modo que a classe de referência é determinada a priori;
para o classificador de movimentos, os participantes executarão cada um dos seis passos
do forró universitário de forma isolada e identificada. As métricas reportadas incluirão
acurácia global, precisão e \textit{recall} por classe, além da matriz de confusão para
identificação de passos sistematicamente confundidos pelo
modelo~\cite{roggio2024pose, bursa2023tflite}.

\subsubsection{Teste com Usuários}
Os testes serão conduzidos mediante aprovação do Comitê de Ética em Pesquisa (CEP) da
UNIFEI e assinatura de Termo de Consentimento Livre e Esclarecido (TCLE) por cada
participante, em conformidade com a Resolução CNS 510/2016. Serão recrutados aproximadamente 10~participantes praticantes de forró universitário
de experiência variada (iniciantes a intermediários) da comunidade de dança da UNIFEI,
tamanho amostral compatível com estudos exploratórios de usabilidade baseados em
escala SUS~\cite{stanford_ai_dance_coaching}. Os participantes
executarão sequências do passo básico com e sem o \textit{feedback} do sistema. A coleta de
dados incluirá: índice de conformidade por módulo --- acurácia dos classificadores de postura
e movimentos, acertos rítmicos e consistência de transferência de peso --- como indicadores
objetivos de desempenho; questionário SUS (\textit{System Usability Scale}) para avaliação de usabilidade
percebida; e entrevista semiestruturada sobre a utilidade e clareza do \textit{feedback}.
Esse protocolo é compatível com a abordagem adotada por Zhou~\cite{stanford_ai_dance_coaching},
que demonstrou, em grupo de dez participantes, ganho de acurácia de 55,9\% para 66,8\% com
\textit{feedback} automatizado baseado em estimativa de pose --- resultado que o presente
experimento buscará verificar no domínio específico do forró universitário.
